\uuid{Z67x}
\titre{Géométrie}
\auteur{bodin ; hanani ; mzari}
\organisation{exo7}
\licence{CC BY-NC-SA 4.0}
\datecreate{2020-01-20}
\niveau{L1}
\module{Algèbre}
\chapitre{Nombres complexes}
\sousChapitre{Racine carrée, équation du second degré}
\difficulte{3}
\theme{Équations, géométrie}
\exo7qcmid{lille-1.nombres-complexes.030}

\contenu{
\begin{question}
\qclassification{104.02,104.03,104.04}
\qtags{motcle=géométrie}
Soit $E$ l'ensemble des points $M$ d'affixe $z$ tels que $M$ et les points $A$ et $B$ d'affixe $i$ et $iz$ respectivement
soient alignés. Quelles sont les assertions vraies ?
\begin{answers}
    \bad{$E$ est la droite passant par les points d'affixe $i$ et $-1+i$ respectivement.}

    \good{$E$ est le cercle de rayon $\frac{1}{\sqrt 2}$  et de centre le point d'affixe $\frac{1}{2}(1+i)$.}

    \bad{$E$ est le cercle de rayon $\frac{1}{2}$  et de centre le point d'affixe $1+i$.}

    \bad{$E$ est la droite passant par les points d'affixe $-i$ et $1-i$ respectivement.}

\end{answers}
\begin{explanations}
$M(z), A(i)$ et $B(iz)$ sont alignés si et seulement si les vecteurs $\overrightarrow{AM}$ et $\overrightarrow{AB}$
sont colinéaires. On pose $z=x+iy, \, x,y\in \Rr$. Les vecteurs $\overrightarrow{AM}$ et $\overrightarrow{AB}$ sont de coordonnées $(x,y-1)$ et $(-y,x-1)$ respectivement. $M(x+iy) \in E$ si et seulement si  $\det(\overrightarrow{AM},   \overrightarrow{AB})=0$.

\end{explanations}
\end{question}
}
